% source: manti.01.ma, 1ère SM, « Fonctions 2 » (Drive 0B6pfSJcEVng0UzNKM1FzeXNsbEU).
% Ex. (x−E(x))/√x = fonctions_s5 Ex.5 ; Ex. a,b,c = fonctions_s4 Ex.3 -> non repris
\ex{Fonctions}{UzNK}{Ex. 1}{6}
Déterminer l'ensemble de définition de la fonction $f$ dans chacun des cas suivants :
\begin{enumerate*}[label=\arabic*),itemjoin={\quad}]
\item $f(x)=\dfrac{x^2+x-2}{|x-4|+3x}$
\item $f(x)=\dfrac{2x^2+x+1}{x|x|+1}$
\item $f(x)=\dfrac{x+2}{2-\sqrt{x^2+x-2}}$
\item $f(x)=\dfrac{x-1}{E(x-2)+1}$
\item $f(x)=\sqrt{E(x)+2}$
\end{enumerate*}
\fin
\ex{Fonctions}{UzNK}{Ex. 2}{8}
Soit $f$ la fonction définie par $f(x)=\sqrt{x-\sqrt{2x-1}}+\sqrt{x+\sqrt{2x-1}}$.
\begin{enumerate}
\item Déterminer $D_f$.
\item Calculer $\big(f(x)\big)^2$.
\item En déduire que $f$ est constante sur l'intervalle $[1,2]$.
\end{enumerate}
\fin
\ex{Fonctions}{UzNK}{Ex. 3}{7}
On considère la fonction $f$ définie par $f(x)=\sqrt{|x|-x}-\sqrt{|x|+x}$.
\begin{enumerate}
\item Déterminer $D_f$ et étudier la parité de $f$.
\item Montrer que pour tout $x$ de $\mathbb{R}^*$ : $f(x)=-x\sqrt{\dfrac{2}{|x|}}$.
\end{enumerate}
\fin
\ex{Fonctions}{UzNK}{Ex. 4}{7}
Soit $f$ la fonction définie par $f(x)=\sqrt{x^2+|x|}-|x|$.
\begin{enumerate}
\item Déterminer l'ensemble de définition de $f$.
\item Montrer que $2|f(x)|<1$ pour tout $x$ de $D_f$.
\end{enumerate}
\fin
